
Beyond the constraints on direct comparison with earlier detectors, this study had several limitations. The benchmark was restricted to both driving corpora, and generalisation to other experimental paradigms or clinical populations was therefore not tested. Event detection was evaluated using a single matching tolerance, although the reported scores would change under a stricter overlap criterion; absolute $F_1$ values should consequently not be compared across studies without accounting for differences in event-matching rules~\citep{salles2024softed}. Recall was not stratified by blink duration, so performance on prolonged eye closures, which are particularly relevant to drowsiness assessment, was not established. Ground truth was based on expert annotation of frontopolar EEG and therefore inherited the annotator's operational definition of blink onset and offset, which remains non-trivial because blink boundaries depend on the signal representation and adopted criteria~\citep{nystrm2024what}. Finally, the proposed method is a thresholding detector and was not evaluated as superior to learned detectors such as transformer-, vision-, and deep-learning-based approaches~\citep{fodor2023blinklinmult,hong2024robust,egambaram2022investigation}. Its intended advantages are instead methodological simplicity, interpretability, and the absence of per-dataset training.
